\chapter{多项式环}\label{chap:polynomial-rings}

\noindent\textbf{第 \thechapter{} 章学习目标}
\begin{itemize}
  \item 定义多项式和多项式环, 并能对它们进行运算.
  \item 定义理想与商并举例, 特别是在多项式环中.
  \item 描述域上多项式的性质, 例如不可约分解的存在性与唯一性.
  \item 判定给定多项式是否不可约.
\end{itemize}

代数几何研究多项式方程的解, 为解集的几何性质与定义这些解集的多项式方程的代数性质建立一套对应关系. 在这个预备章节中, 我们将介绍学习代数几何所需的多项式环的代数概念. 本章最后将证明一个重要事实: 域上的每个多项式都能唯一地分解成不可约多项式的乘积.

我们假定读者已经修过一门环论入门课程, 熟悉环, 整环和域的概念. 不过, 并不要求读者掌握最基本的定义和结果以外的知识. 代数基础较扎实的学生自然能更快读完本章, 而刚接触抽象代数的学生则可以选择花较多时间掌握这几页的内容.

本章旨在为全书其余部分建立代数基础, 着重讨论对后续内容最有用的多项式环基本性质. 因此, 为求简明, 我们常常不追求一般性. 例如, 我们不会考虑非交换环, 所以每个理想都是双边理想. 本章几乎所有例子都在多项式环的具体情形中说明这些概念, 不过我们希望在这一背景下培养的直觉, 能帮助有兴趣的学生在更一般的情形中研究这些概念.

\par\smallskip
{\footnotesize
\noindent\textcopyright{} The Author(s) 2025\par
\noindent E. Clader, D. Ross, \emph{Beginning in Algebraic Geometry}, Undergraduate Texts in Mathematics,
\href{https://doi.org/10.1007/978-3-031-88819-9_1}{\nolinkurl{https://doi.org/10.1007/978-3-031-88819-9_1}}.\par}
